\documentclass[11pt,letterpaper]{article}

%====================================================================
% 1. MATEMÁTICA Y FUENTES
%====================================================================

% Compilador: XeLaTeX
\usepackage{fontspec}
\usepackage{amsmath}
\usepackage{mathtools}
\usepackage{amsthm}
\usepackage{amssymb}
\usepackage{unicode-math}

%====================================================================
% 2. IDIOMA
%====================================================================

\usepackage{polyglossia}
\setmainlanguage{spanish}

\usepackage{csquotes}
\usepackage{etoolbox}

%====================================================================
% 3. DISEÑO Y FORMATO GENERAL
%====================================================================

\usepackage[margin=1in]{geometry}

\usepackage{microtype}
\usepackage{setspace}
\usepackage{enumitem}
\usepackage{ragged2e}

% Interlineado
\setstretch{1.1}

% Configuración de párrafos
\setlength{\parindent}{1.2em}
\setlength{\parskip}{0pt}
\setlength{\emergencystretch}{2em}

% Configuración de listas
\setlist{
    itemsep=2pt,
    topsep=4pt
}

%====================================================================
% 4. TABLAS
%====================================================================

\usepackage[table]{xcolor}

\usepackage{booktabs}
\usepackage{tabularx}
\usepackage{longtable}
\usepackage{array}

%====================================================================
% 5. FIGURAS Y ELEMENTOS GRÁFICOS
%====================================================================

\usepackage{graphicx}
\usepackage{float}
\usepackage{caption}
\usepackage{pdflscape}

% Configuración de títulos de tablas y figuras
\captionsetup{
    font=small,
    labelfont=bf,
    labelsep=period,
    skip=6pt
}

%====================================================================
% 6. CÓDIGO Y LISTADOS
%====================================================================

\usepackage{listings}

\lstset{
    language=Python,
    basicstyle=\ttfamily\small,
    frame=lines,
    keywordstyle=\bfseries,
    commentstyle=\itshape\color{gray},
    showstringspaces=false,
    columns=fullflexible,
    xleftmargin=1em
}

%====================================================================
% 7. ENCABEZADOS Y PIES DE PÁGINA
%====================================================================

\usepackage{fancyhdr}

\setlength{\headheight}{15pt}

\pagestyle{fancy}
\fancyhf{}

\renewcommand{\headrulewidth}{0pt}

\fancyhead[R]{
    \small\itshape Informe de avance del proyecto
}

\fancyfoot[C]{
    \small\thepage
}

%====================================================================
% 8. BIBLIOGRAFÍA
%====================================================================

% Procesamiento mediante Biber
\usepackage[
    backend=biber,
    style=authoryear,
    sorting=nyt
]{biblatex}

% Archivo de referencias bibliográficas
\addbibresource{references.bib}

%====================================================================
% 9. ENLACES Y REFERENCIAS CRUZADAS
%====================================================================

\usepackage{xurl}

\usepackage[
    colorlinks=true,
    linkcolor=blue!40!black,
    citecolor=blue!40!black,
    urlcolor=blue!40!black,
    pdfusetitle
]{hyperref}

\hypersetup{
    pdftitle={Informe de avance del proyecto},
    pdfauthor={Anthonny Flores Rojas,
               Andrey Gonzalez Bastos,
               Diego Vega}
}

\usepackage[capitalise]{cleveref}

% Nombres utilizados en las referencias cruzadas
\crefname{theorem}{Teorema}{Teoremas}
\crefname{proposition}{Proposición}{Proposiciones}
\crefname{lemma}{Lema}{Lemas}
\crefname{corollary}{Corolario}{Corolarios}
\crefname{definition}{Definición}{Definiciones}
\crefname{example}{Ejemplo}{Ejemplos}
\crefname{conjecture}{Conjetura}{Conjeturas}
\crefname{section}{Sección}{Secciones}
\crefname{table}{Tabla}{Tablas}
\crefname{figure}{Figura}{Figuras}
\crefname{equation}{ecuación}{ecuaciones}

%====================================================================
% 10. TEOREMAS Y ENTORNOS MATEMÁTICOS
%====================================================================

% Numeración de teoremas por sección
\theoremstyle{plain}

\newtheorem{theorem}{Teorema}[section]
\newtheorem{proposition}[theorem]{Proposición}
\newtheorem{lemma}[theorem]{Lema}
\newtheorem{corollary}[theorem]{Corolario}
\newtheorem{conjecture}[theorem]{Conjetura}

\theoremstyle{definition}

\newtheorem{definition}[theorem]{Definición}
\newtheorem{example}[theorem]{Ejemplo}

\theoremstyle{remark}

\newtheorem{remark}[theorem]{Observación}

% Numeración de ecuaciones por sección
\numberwithin{equation}{section}

% Nombre de las demostraciones
\gappto\captionsspanish{%
    \renewcommand{\proofname}{Demostración}
}

%====================================================================
% 11. COMANDOS MATEMÁTICOS
%====================================================================

\newcommand{\R}{\mathbb{R}}
\newcommand{\Z}{\mathbb{Z}}
\newcommand{\E}{\mathbb{E}}

\renewcommand{\P}{\mathbb{P}}

\newcommand{\dd}{\,\mathrm{d}}

\newcommand{\norm}[1]{%
    \lVert #1\rVert
}

\newcommand{\Sh}{\mathcal{S}}

\DeclareMathOperator{\LSE}{LSE}

%====================================================================
% 12. ENTORNOS MATEMÁTICOS EN ESPAÑOL
%====================================================================

\newtheorem{teorema}[theorem]{Teorema}
\newtheorem{proposicion}[theorem]{Proposición}
\newtheorem{lema}[theorem]{Lema}
\newtheorem{corolario}[theorem]{Corolario}
\newtheorem{definicion}[theorem]{Definición}
\newtheorem{ejemplo}[theorem]{Ejemplo}
\newtheorem{observacion}[theorem]{Observación}

%====================================================================
% 13. COMANDOS PARA MONEDAS
%====================================================================

\newcommand{\eur}{\texttt{EUR}\,}
\newcommand{\gbp}{\texttt{GBP}\,}
\newcommand{\czk}{\texttt{CZK}\,}

%====================================================================
% 14. CLASIFICACIÓN DE DATOS
%====================================================================

% Categoría A: datos publicados
\newcommand{\catA}{\textbf{[A]}}

% Categoría B: cálculos propios
\newcommand{\catB}{\textbf{[B]}}

% Categoría C: datos no disponibles o no verificables
\newcommand{\catC}{\textbf{[C]}}

%====================================================================
% 15. COMANDOS AUXILIARES DEL INFORME
%====================================================================

% Resaltar el estado de un elemento
\newcommand{\estado}[1]{\textbf{#1}}

% Entornos de interlineado sencillo
\newcommand{\sw}{\begin{singlespace}}
\newcommand{\ew}{\end{singlespace}}

% Texto provisional
\newcommand{\ph}[1]{%
    \textcolor{red!70!black}{\textit{[#1]}}
}

% Estados para tablas de avance
\newcommand{\done}{%
    \textcolor{green!45!black}{\textbf{Completado}}
}

\newcommand{\wip}{%
    \textcolor{orange!80!black}{\textbf{En curso}}
}

\newcommand{\todo}{%
    \textcolor{red!70!black}{\textbf{Pendiente}}
}

%====================================================================
% 16. DATOS INSTITUCIONALES
%====================================================================

\newcommand{\Universidad}{Universidad de Costa Rica}
\newcommand{\Escuela}{Escuela de Matemáticas}

\newcommand{\Curso}{
    Valoración de instrumentos financieros (CA0308)
}

\newcommand{\Profesor}{
    Daniel Vargas Carrillo
}

%====================================================================
% 17. DATOS DEL DOCUMENTO
%====================================================================

\title{
    \textbf{Nombre del informe}\\[0.4em]
    \large Informe de anteproyecto
}

\author{%
    {Anthonny Flores Rojas}
    \thanks{
        Escuela de Matemáticas,
        Universidad de Costa Rica.
        \texttt{anthonny.flores@ucr.ac.cr}
    }
    \and
    {Andrey Gonzalez Bastos}
    \thanks{
        Escuela de Matemáticas,
        Universidad de Costa Rica.
        \texttt{andrey.gonzalezbastos@ucr.ac.cr}
    }
    \and
    {Diego Vega}
    \thanks{
        Escuela de Matemáticas,
        Universidad de Costa Rica.
        \texttt{diego.vega@ucr.ac.cr}
    }
}

\date{
    Curso: Valoración de instrumentos financieros (CA0308)\\
    Profesor: Daniel Vargas Carrillo\\[0.3em]
    \today
}

\begin{document}
	
	
	\maketitle
	\tableofcontents   % opcional; muchos papers largos lo incluyen
	
% =====================================================================
\section{Introducción}\label{sec:intro}

El sector de la industria energética es quizá uno de los más importantes en la economía mundial, gracias a él, el mundo peude movilizarse gracias al petróleo que termina en gasolina, o comer, gracias al gas que permite cocinar alimentos. El negocio de este sector es el de proveer todas las materias primas que se requieren para energía, tales como el petróleo, gas o electricidad. En el caso de la presente investigación, Giannozzi y Kirilova (2026) plantean el estudio de un conjunto de 120 empresas dedicadas exclusivamente a este sector. \vspace{0.50 cm}

Las empresas a nivel mundial siempre buscan mitigar o reducir al mínimo su exposición al riesgo, esto mediante herramientas financieras como los derivados, los cuales siempre están ligados a un activo subyacente. Sin embargo, esta estrategia también conlleva sus riesgos y debe ser analizada en profundidad antes de ser presentada y ejecutada. En ciertos contextos, puede más bien magnificar de forma imprevista el riesgo al que se enfrentan las empresas, sobre todo en tiempos en los que los mercados pasan por un fuerte estrés o alta volatilidad. 
\vspace{0.50 cm}

En el artículo, como las empresas se cubren con futuros, obviamente están expuestos a que los activos subyacentes ligados a estos sufran variaciones y pierdan o ganen valor, obviamente si pierden valor la empresa se expone a grandes pérdidas.
\vspace{0.50 cm}

Un caso específico de exposición es el margin call, el cual, según Interactive Brokers (2026) se produce cuando un corredor exige a un inversor que deposite más dinero o valores en su cuenta de margen para alcanzar el mínimo de margen requerido. Esto suele ocurrir cuando el valor de los valores en la cuenta de margen cae por debajo de un determinado umbral, lo que provoca que el capital de la cuenta disminuya por debajo del margen exigido.
\vspace{0.50 cm}

El margin call depende de la posición que se tenga, no del movimiento del subyacente como tal. El margin call se da cuando el precio del subyacente se mueve en contra de la posición, es decir:
\vspace{0.50 cm}

\begin{itemize}
    \item \textbf{Posición corta:} El margin call llega cuando el subyacente sube.
    \item \textbf{Posición larga:} El margin call llega cuando el subyacente baja.
\end{itemize}
\vspace{0.50 cm}

Para el estudio del artículo, se toma en cuenta cuando el subyacente sube, esto dado que cubren el riesgo de que los precios de sus producciones bajen. Es decir, pactan un precio a futuro sobre su producción con \textbf{posiciones cortas} respecto a sus derivados. Acá puede existir confusión ya que la empresa se encuentra en una \textbf{posición larga} respecto al activo físico como tal, a su producción por ejemplo de gas natural, pero, a la hora de hablar del derivado, ellos están cortos porque pactan el precio a futuro para cubrirse del riesgo de que los precios se desplomen en el mercado, más aún tomando en cuenta que los derivados negociados en bolsa se valoran mark-to-market, es decir, día a día, aumentando mucho la exposición. 
\vspace{0.50 cm}

Los autores Giannozzi y Kirilova (2026) dicen ser los primeros en descubir que, aunque el dinero de esos margin call se devuelve cuando los precios se normalizan, el efecto no es transitorio como se pensaba, sino que las consecuencias en el riesgo de default, el costo de deuda y en los márgenes de utilidad pueden persistir por más de un año luego del llamado. 
\vspace{0.50 cm}

En el caso de estudio, se investiga sobre la crisis energética ocurrida en el 2022 en Europa. El artículo explica que la decisión de tomar este episodio se fundamenta en que la industria energética usa mucho los derivados como su estrategia de cobertura ante riesgos y tienen un alto potencial de impacto en la salud financiera de la compañía. Por otro lado, también se toma en cuenta este sector y episodio dado que los requerimientos de margen son sumamente altos comparados con otros derivados, por ejemplo, el gas y el petróleo ven fluctuaciones del 75\% y 41\% respectivamente, mientras que el par EUR/USD sufrió una fluctuación de apenas el 8\% durante el mismo período, esto hace que un llamado de margen pueda ser devastador para esas empresas. Por último, se justifica la elección porque se pueden apreciar mejor las consecuencias del riesgo amplificado y tiene una relación estrecha con la política. 
\vspace{0.50 cm}

La fecha de corte del caso de estudio es el año 2022, año marcado por la política y la guerra entre Rusia y Ucrania. Los precios que interesan son los del petróleo, la elecricidad y el gas. Para el petróleo, interesa el futuro de Brent Crude Oil, según Yahoo Finance, el año empezó en \$81.75, llegando a un máximo de \$122.01, por lo que podemos ver las altas fluctuaciones que se viven en estos mercados. Otro ejemplo muy claro sucede con el futuro Dutch TTF Natural Gas Calendar, el cual arrancó el 2022 con un precio de \$88.174 y llegando a un máximo de \$339.196. Por último, a inicios del 2022, una corona checa equivalía a aproximadamente 0.04572 según la plataforma Wise.
\vspace{0.50 cm}

Como se puede apreciar en el párrafo anterior, los precios de los subyacentes se dispararon durante el año más fuerte del conflicto entre Rusia y Ucrania. Entonces, como se había dicho, como tienen posiciones cortas en el derivado, estas alzas y estrés en los mercados, provocaron que las empresas tuvieran fuertes margin calls durante ese período, haciendo así que los efectos negativos se magnificaran dado el fuerte disparo en los precios de los subyacentes. Esto se suma a que si la empresa no pagase el margin call, la contraparte puede cerrar las posiciones, realizando la pérdida antes del vencimiento, eso expone a la empresa a un estrés financiero fuerte o hasta una eventual quiebra. 


\vspace{0.50 cm}

Como el caso de estudio es sumamente grande (120 empresas), por lo cual se decide para el análisis del presente trabajo tomar únicamente en cuenta al grupo empresarial \v{C}EZ. Esta empresa radica en  Praga, República Checa. Su negocio principal es el de generar, distribuir, comerciar y vender electricidad y calor, así como gas natural. La empresa, según Prague Stock Exchange, empezó a cotizar en bolsa de valores en 1993, a la fecha de esta investigación, la empresa tiene una cotización en bolsa de valores de 1366 coronas checas, lo que equivale a aproximadamente 63 dólares estadounidenses. A inicios del 2022, tenía una cotización de 811 coronas checas, por lo que vemos que en este período de tiempo su valor ha ido en aumento. Es importante señalar que la empresa cuenta con entrega de dividendos a sus accionistas, el último fue de 42 coronas checas, de aproximadamente unos 2 dólares estadounidenses. 
\vspace{0.50 cm}

En conclusión, cuando los mercados presentan grandes fluctuaciones, se suele hacer este margin call que puede afectar directamente las utilidades de la empresa y representar un dolor de cabeza mayor que el riesgo original que se pretendía mitigar. El caso habla de que se estima que la industria Europea de energía necesitó más de 1.5 trillones de euros para cubrirse del costo de los llamados de margen. 

	\subsection{Selección de la empresa y delimitación del análisis}

	\noindent El estudio se centra en ČEZ, a.~s., grupo energético europeo dedicado principalmente a la generación y comercialización de electricidad. Cuando se utilizan cifras de la empresa, se emplea información consolidada del grupo, sin confundir los datos de sus filiales con los del conjunto.

	\vspace{0.2cm}

	\noindent Dentro de esta actividad, 
	el análisis se limita a la generación de electricidad y a 
	las ventas a plazo mediante las cuales el grupo 
	cubre parte de su producción futura \parencite{cez_q3_2021}. 
	La variable principal de interés es el precio de la electricidad. 
	La documentación disponible no permite todavía atribuir a la empresa un contrato bursátil específico, por lo que no se asigna ninguno en esta etapa.

	\vspace{0.2cm}

	\noindent Un generador mantiene una exposición económica larga 
	al precio de la electricidad, porque sus ingresos aumentan cuando 
	este sube. Al vender por anticipado parte de su producción a un 
	precio fijado, adopta además una posición corta en el mercado a 
	plazo. Si el precio aumenta, esa posición pierde valor y, cuando 
	el contrato está sujeto a garantías, obliga a aportar margen 
	adicional \parencite{GIANNOZZI2026107700}. El alcance exacto de 
	los derivados de ČEZ sometidos a garantías requiere verificación 
	documental. Por ello, el análisis se concentra en el riesgo de 
	liquidez: un aumento extraordinario del precio de la electricidad, 
	que constituye un riesgo de mercado, se traduce en llamadas de margen que 
	consumen efectivo antes de que se perciban los ingresos de la producción cubierta.
	

	\newpage

\printbibliography[title=Referencias]

\end{document}